\documentclass{sajzk}

\usepackage{array}

\usepackage{minted}
% \usepackage{listings}
% \lstset{basicstyle=\ttfamily}
% \usepackage{xcolor}
% \usepackage{mdframed}
% 
% \definecolor{light-gray}{gray}{0.95} % the shade of grey that stack exchange uses

\usetikzlibrary {arrows.meta, angles, quotes}
\usetikzlibrary {calc}

\begin{document}
\section{Vergleich der Beweismöglichkeiten} 
\label{ed3f.tex}
\lhead{:mathe:logik:kalkül:}

In dem Buch von Wolfgang Seigmüller werden für die Beweisfürungen drei
Möglichkeiten beschrieben

\begin{itemize}
  \item Baumkalkül
  \item Gentzen-Kalkül
  \item Hilbert-Kalkül
  \item Schütte-Kalkül
\end{itemize}

Es werden auch die Equivalienzen der einzelne Kalküle betrachtet.

\subsection{Referenzen} 
\begin{itemize}
  \item Beweismöglichkeiten Baumkalkül (asdf.md)
\end{itemize}

\subsection{Literatur} 
\begin{itemize}
  \item Wolfgang Stegmüller, Matthias Varga von Kibéd - Strukturtypen der Logik
\end{itemize}
\end{document}
